% --------------------------------
% RESEARCH
% --------------------------------


\mediumtitle{Research posts}

\begin{doubletablelist}
	
	\dtlist{Feb 2025 - Present}{\textbf{Visiting Senior Lecturer} at the Department of Twin Research \& Genetic Epidemiology, King's College London \\	
	 &\textsc{Research activity}: Alessia Visconti continues her research in \emph{(a)}~IgA glycosylation~\cite{Vis24}, \emph{(b)}~genetic epidemiology~\cite{Ros24,Sta25,Ali26}, and \emph{(d)}~the gut microbiome and its effects on human health~\cite{Dan24,Att24,Val24,Sar25,Pop25,Lin26}, with a particular interest in its phage component~\cite{Kir24,Zol24}. \\
	 &\textsc{Note}: Alessia Visconti was \emph{Visiting Research Fellow} at the same institution from October 2023 to September 2024.
	}

	\dtlist{Dec 2023 - Present}{\textbf{Assistant professor (tenure track)} at the Center for Biostatistics, Epidemiology and Public Health, Department of Clinical and Biological Sciences, University of Turin \\
	 &\textsc{Research activity}: Alessia Visconti's research focuses on \emph{(a)}~the development and application of statistical and computational multi\emph{-omics} approaches to understand the underlying mechanisms, and to improve the diagnosis and treatment of complex human diseases, \emph{(b)}~the design and deployment of deep learning approaches for monitoring and improving public health responses~\cite{Kou25,DeL25,Kou26,Con24}, and \emph{(c)}~the integration of patient preferences into clinical care and drugs and medical devices development~\cite{DiB24}. She continues to teach undergraduate and postgraduate courses, supervise PhD students and postdoctoral researchers, and serve as MSc thesis supervisor. She contributes to the design and writing of international funding proposals, for which she has scientific responsibility. 
	}

	\dtlist{Oct 2023 - Nov 2023}{\textbf{Senior bioinformatician} at the Genomics Research Centre, Human Technopole, Milan \\
	 &\textsc{Research activity}: Alessia Visconti was involved in a project conducting a genome-wide association study (GWAS) on phenotypes extracted from cardiac MRI scans using an unsupervised deep learning approach~\cite{Ome24}.
	}

	\dtlist{Aug 2022 - Apr 2023}{\textbf{Career break}: Maternity leave 
	}

	\dtlist{Aug 2017 - Sep 2023}{\textbf{Research fellow} at the Department of Twin Research \& Genetic Epidemiology, King's College London \\
	 &\textsc{Research activity}: Alessia Visconti's research focused on the development and application of statistical and computational multi\emph{-omics} approaches to understand the underlying mechanisms, and to improve the diagnosis and treatment of complex human diseases. She continued to teach postgraduate courses, supervise PhD students and postdoctoral researchers and serve as MSc thesis supervisor. She contributed to the design and writing of international funding proposals, for which she had scientific responsibility. \\
	 & Building on her previous position, she collaborated on projects studying IgA nephropathy~\cite{Dot21} as well as melanoma and its risk factors, particularly from a genomic perspective~\cite{Vis19a,Vis20,Lan20,San20}, and begun research on predicting the response to immunotherapy in melanoma patients~\cite{Ros22,Vis23}. She also studied \emph{(a)}~the influence of the gut microbiome on human health~\cite{Vis19,Bar20,LeR22,Zha22,Val23,Lou23,Nog23a,Nog23b}, \emph{(b)}~the multi\emph{-omics} profiles underlying chronic autoimmune thyroiditis~\cite{Mar20}, \emph{(c)}~the genomic basis of atopic dermatitis~\cite{Gro21,Bud22}, \emph{(d)}~the influence of amylase copy-number variation (CNV) on a range of adiposity traits~\cite{Ros21}, \emph{(e)}~the effect of cigarette smoking on more than 35,000 immunophenotypes~\cite{Pia21}, and \emph{(f)}~the effect of X-chromosome inactivation~\cite{Zit23}. During the SARS-CoV-2 pandemic, she collaborated within the ZOE COVID Study (now ZOE Health Study) to develop the pipeline for analysing data submitted daily by approximately 4 million users~\cite{Mur21}. She also led two studies on the cutaneous symptoms of the infection~\cite{Vis21, Vis22}, and collaborated as data analyst in several other studies~\cite{Men20,Lee20,Zaz20,Hop21,Wil21,Sud21}. 
	% &\textsc{Supervisor}: Dr. Mario Falchi
	}
%
% 	%     \dtlist{Jun 2016 - Jun 2019}{\textbf{Honorary research associate} at CERN OpenLab, Switzerland \\
% 	%  &\textsc{Research activity}: Alessia Visconti extended the ROOT library (\url{https://root.cern/}), developed by the CERN OpenLab for enabling the storage and scientific analyses and visualisation of large amounts of data from particle physics experiments, to allow the efficient storage of genomic data. \\
% 	%  &
% 	% % &\textsc{Supervisor}: Dr. Marco Manca
% 	% }
%
%
	\dtlist{Apr 2015 - Jul 2017}{\textbf{Research associate} at the Department of Twin Research \& Genetic Epidemiology, King's College London \\
	 &\textsc{Research activity}: Alessia Visconti's research focused on the development and application of statistical and computational multi\emph{-omics} approaches to understand the underlying mechanisms, and to improve the diagnosis and treatment of complex human diseases. In particular, she led the statistical analysis in a series of projects aimed at studying IgA Nephropathy~\cite{Lom16}, cognition and neurodevelopmental disorders~\cite{Cul18}, and epigenetic modification~\cite{Zag18} as well as at dissecting the aetiology of melanoma, melanoma risk phenotypes, and their connection with ageing~\cite{Rib16,Hys18,Vis18a,Duf17}. Becoming interested in the human microbiome, she defined a set of guidelines for conducting metagenomic studies in clinical settings~\cite{Vis18c}, developed a reproducible pipeline for the analysis of metagenomic data (\emph{YAMP})~\cite{Vis18b}, and became the scientific lead for the analysis of metagenomic data at the Department of Twin Research \& Genetic Epidemiology, overseeing the processing and management of more than 4,000 samples (an activity that is still ongoing as part of Alessia Visconti's position as Visiting Senior Lecturer in the Department). \\
	 & She begun to teach postgraduate courses, supervise PhD students and postdoctoral researchers and serve as MSC thesis supervisors. She contributed to the design and writing of international funding proposals, for which she had scientific responsibility. 
	% &\textsc{Supervisor}: Dr. Mario Falchi
	}
%
    \dtlist{Jan 2014 - Mar 2015}{\textbf{Research associate} at the Department of Genomics of Common Disease, School of Public Health, Imperial College London  \\
	 &\textsc{Research activity}: Alessia Visconti developed and implemented an approach for the association testing of quantitative traits on related individuals (\emph{PopPAnTe})~\cite{Vis16}, which was used in one of the first epigenome-wide association studies (EWAS) in a Middle Eastern population~\cite{AlM15}. She collaborated on studies investigating the effect of CNV on reading abilities~\cite{Gia16}, and the gene regulatory networks underlying cognitive abilities associated with neurodevelopmental disorders~\cite{Joh15}. 
	% &\textsc{Supervisor}: Dr. Mario Falchi
	}



	\dtlist{Jan 2012 - Dec 2013}{\textbf{Research associate} at the Department of Computer Science, University of Turin \\ %Recipient of a 2-year fellowship for her research proposal, awarded by the Italian Minister of Education, University and Research. \\
	 &\textsc{Research activity}: Alessia Visconti \emph{(a)}~developed and implemented a bi-clustering approach which can leverage \emph{a priori} knowledge to identify clinically relevant groups (\emph{AID-ISA})~\cite{Vis13a}, \emph{(b)}~contributed to the development of an exact algorithm for answering Maximum a Posteriori queries on tree-structured probabilistic graphical models (\emph{CDoT})~\cite{Esp13}, \emph{(c)}~applied machine learning approaches for predicting small peptide lipophilicity~\cite{Vis15a}, \emph{(d)}~conducted the bioinformatics analysis to assessed the ability of more than 200 compounds to act as hydrogen-bond donors~\cite{Erm14}, and \emph{(e)}~developed a tool for searching topologies within chemical molecules, which has been used to support teaching at the Department of Molecular Biotechnology and Health Sciences, University of Turin (\emph{PicNic}).
	% &\textsc{Supervisors}: Dr. Roberto Esposito
	}


	\dtlist{Jun 2011 - Dec 2011}{\textbf{Visiting researcher} at the Center of Biological Sequence Analysis, Technical University of Denmark
	\\ % at the Regulatory Genomics Group at the Center of Biological Sequence Analysis, Systems Biology Department, Technical University of Denmark, Denmark
	 &\textsc{Research activity}: Alessia Visconti used data generated at the Danish centre to infer causal relationships between transcriptomic and proteomic levels using Granger causality testing~\cite{Vis12b}.  She developed, together with Ali Altıntas and Chris Workman, an \emph{ensemble} approach combining results from regression trees and support vector machine for regression for the \emph{``DREAM6 -- Promoter Activity Prediction Challenge''}, which ranked 8th out of 21 participants.
	% &\textsc{Supervisor}: Prof. Christopher Workman
	}


	\dtlist{Jan 2009 - Dec 2011}{\textbf{PhD student} at Department of Computer Science, University of Turin \\
	  &\textsc{Research activity}: Alessia Visconti developed and implemented: \emph{(a)}~an algorithm for extracting protein motifs and clustering them into families using a constrained co-clustering approach (also providing a web interface, \emph{MotifsLinker})~\cite{Cor09a}, \emph{(b)}~a co-clustering algorithm embedding Gene Ontology annotations to provide more meaningful grouping (again accompanied by a web interface, \emph{GOClust})~\cite{Vis11c,Cor09b}, \emph{(c)}~an algorithm for clustering Gene Ontology terms according to their function to obtain a more compact and informative ontology (\emph{RGO})~\cite{Vis11a,Vis10a}, \emph{(d)}~two algorithms for the reverse engineering of gene regulatory networks~\cite{Vis11b,Mar12,Vis12b}, and \emph{(e)}~contributed to the development of a \\
	  & modular software architecture for the analysis of metagenomic sequences, taking primary responsibility for the co-clustering module~\cite{Bon11}.
	  	% &\textsc{Supervisors}: Dr. Roberto Esposito, Prof. Marco Botta
	}


% 	\dtlist{Sep 2008 - Dec 2008}{\textbf{Research assistant} at the Department of Computer Science, University of Turin and in collaboration with the Department of Arboriculture and Pomology, University of Turin, Italy  \\
%  	 &\textsc{Research activity}: During this post, Alessia Visconti contributed to the development of computational approaches for the classification and traceability of fruits produced in Piemonte. \\
% 	 &
% % 	% &\textsc{Supervisors}: Prof. Marco Botta, Prof. Roberto Botta
%  	}

 \end{doubletablelist}
 
 
 
 % --------------------------------
 % RESEARCH VISIT

 % \dtlist{Research Visits}{
 % 	\vspace{-0.6cm}
 % 	\begin{itemize} %[itemsep=-0.5ex]
 % 		\minusitem \begin{minipage}{0.65\textwidth}
 % 		 \textsc{Jan 2015} -- visiting researcher at Weill Cornell Medical College in Qatar
 % 		\end{minipage}
 % 		\minusitem \begin{minipage}{0.65\textwidth}
 % 		 \textsc{Jun - Dec 2011} -- visiting PhD student at the Regulatory Genomics Group, Center of Biological Sequence Analysis, Systems Biology Department, Technical University of Denmark, under the supervision of Prof. C. Workman
 % 		\end{minipage}
 % 	\end{itemize}
 % }